% !PRINTED 329
\weberpage{329}

But if we now let $x$ approach the root $\alpha$, then $X$ will approach the same root, and (20) shows that $\xi$ also tends to the same limit. It follows, therefore, that if $x$ lies sufficiently near a root, this nearest root lies in the interval $(\xi, X)$, in which no second root is contained, and that this interval closes more and more about $\alpha$.

As an example we take (with \gesperrt{Laguerre}) the spherical functions, of which we have already shown earlier (\S.\,87) that they have only real roots, which all lie between $-1$ and $+1$. The function
\[
P_n(x) = \frac{1\,.\,3 \ldots (2\,n-1)}{1\,.\,2 \ldots n} \left( x^n - \frac{n(n-1)}{2(2\,n-1)}\,x^{n-2} + \cdots \right)
\]
has the property that
\[
P_n(1) = 1, \qquad P_n'(1) = \frac{n(n+1)}{2},
\]
\[
P_n''(1) = \frac{n(n-1)(n+1)(n+2)}{8}
\]
holds. One proves this easily by complete induction, by verifying the formulas directly in the first cases $n = 2, 3 \ldots$, and then, from the recursion formula
\[
(1 - x^2)\,P_n'(x) + n\,x\,P_n(x) - n\,P_{n-1}(x) = 0
\]
by differentiating twice, establishing their correctness for the index $n$, provided it is assumed for the index $n-1$.

If we then set $x = 1$ in (19) and (20), we obtain two limits between which the root coming nearest to 1 lies, namely
\[
X = 1 - \frac{2}{n + 1 + (n-1)\sqrt{\dfrac{n(n+1)}{2}}},
\]
\[
\xi = 1 - \frac{2}{n + 1 + \sqrt{\dfrac{n(n+1)}{2}}}.
\]
For $n = 7$, for example, this gives
\[
X = 0{,}94967 \ldots
\]
\[
\xi = 0{,}84952 \ldots
\]
